\section{A recursive formula for plethysm coefficients}\label{sec:14}
The main result of this article is Theorem A, a recursive formula for plethysm coefficients of the form $a^\mu_{\lambda,(m)}$ for arbitrary $m\in\N$ and partitions $\mu$ and $\lambda$. Together with Lemma~\ref{lem:tall-pleth}, it describes the deflations $\delta^\mu$ of $\mu\vdash mn$ with respect to $S_n$, noting that $a^\mu_{\lambda',(m)}=\langle \sgn_{S_n}\cdot\ \delta^\mu,\chi^\lambda\rangle$ for $\lambda\vdash n$. We restate Theorem A here as Theorem~\ref{thm:14.6i} for ease of reference for the reader, and recall that plethysm coefficients indexed by skew shapes were defined in \eqref{eqn:skew-plethysm}:

\begin{theorem}[Theorem A]\label{thm:14.6i}
	Fix $n\in\N$. Let $m\in\N$, $k\in\{0,1,\dotsc,n-1\}$ and $\lambda\vdash n$. Let $\mu\vdash mn$ with $l(\mu)=n-k$, and set $\hat{\mu}\coloneqq \mu-(1^{n-k})\vdash (m-1)n+k$. Then
	\begin{equation}\label{eqn:A}
		a^\mu_{\lambda',(m)} = \sum_{i=0}^k (-1)^{k+i} \cdot \sum_{\substack{\alpha\vdash k+(m-1)i\\\beta\vdash i}} a^{\alpha/(k-i)}_{\beta',(m)}\cdot a^{\hat{\mu}/\alpha}_{\lambda/\beta,(m-1)}.
	\end{equation}
\end{theorem}

We first illustrate some of the uses of our main theorem in Example~\ref{ex:14.7} and in proving a stability result (Proposition~\ref{prop:17}), before proving Theorem~\ref{thm:14.6i} in full in Section~\ref{sec:16}. We present further applications to Sylow branching coefficients in Section~\ref{sec:13}.

\begin{example}\label{ex:14.7}
	We illustrate how to compute $a^{(6,3^{n-2})}_{\lambda,(3)}$ for all $n\ge 6$ and $\lambda\vdash n$ using Theorem~\ref{thm:14.6i}. 
	Let $\mu=(6,3^{n-2})$. Since $l(\mu)=n-1$, Theorem~\ref{thm:14.6i} gives
	\begin{align*} a^\mu_{\lambda,(3)}& = - a^{\hat{\mu}/(1)}_{\lambda',(2)} + a^{\hat{\mu}/(3)}_{\lambda'/(1),(2)} \\&= -a^{(4,2^{n-2})}_{\lambda',(2)} - a^{(5,2^{n-3},1)}_{\lambda',(2)} + a^{(2^{n-1})}_{\lambda'/(1),(2)} + a^{(3,2^{n-3},1)}_{\lambda'/(1),(2)} + a^{(4,2^{n-3})}_{\lambda'/(1),(2)}, \end{align*}
	since $a^\emptyset_{\emptyset,(3)}=1$ and $a^\alpha_{(1),(3)}=\delta_{\alpha,(3)}$. Applying Theorems~\ref{thm:14.6i} and~\ref{thm:LR}, we obtain
	\begin{align*}
		a^{(4,2^{n-2})}_{\lambda',(2)} &= \begin{cases} 1 & \text{if }\lambda'\in\{(n),(n-1,1),(n-2,2)\},\\ 0&\text{otherwise},\end{cases}\\
		a^{(5,2^{n-3},1)}_{\lambda',(2)} &= \begin{cases} 1 & \text{if }\lambda'\in\{(n-1,1),(n-2,2),(n-2,1^2),(n-3,2,1)\},\\ 0&\text{otherwise}.\end{cases}
	\end{align*}
	Similarly, since $a^\theta_{\lambda'/(1),(2)} = \sum_{\varepsilon\vdash n-1}c^{\lambda'}_{\varepsilon,(1)}\cdot a^\theta_{\varepsilon,(2)}$, we have that
	\begin{align*}
		a^{(2^{n-1})}_{\lambda'/(1),(2)} &= c^{\lambda'}_{(n-1),(1)} = \begin{cases}
			1 & \text{if }\lambda'\in\{(n),(n-1,1)\},\\ 0&\text{otherwise},\end{cases}\\
		a^{(3,2^{n-3},1)}_{\lambda'/(1),(2)} &= c^{\lambda'}_{(n-2,1),(1)} = \begin{cases}
			1 & \text{if }\lambda'\in\{(n-1,1),(n-2,2),(n-2,1^2)\},\\ 0&\text{otherwise},\end{cases}\\
		a^{(4,2^{n-3})}_{\lambda'/(1),(2)} &= c^{\lambda'}_{(n-2,1),(1)} + c^{\lambda'}_{(n-1),(1)} + c^{\lambda'}_{(n-3,2),(1)}\\
		&= \begin{cases}
			2 & \text{if }\lambda'\in\{(n-1,1),(n-2,2)\},\\
			1 & \text{if }\lambda'\in\{(n),(n-2,1^2),(n-3,3),(n-3,2,1)\},\\
			0 & \text{otherwise}.
		\end{cases}
	\end{align*}
	Putting this together, we obtain
	\[ a^\mu_{\lambda,(3)} = \begin{cases}
		2 & \text{if }\lambda'=(n-1,1),\\
		1 & \text{if }\lambda'\in\{(n),(n-2,2),(n-2,1^2),(n-3,3)\},\\
		0 & \text{otherwise}.
	\end{cases} \]
	In particular, this gives an alternative method for one of the steps in the proof of Theorem~\ref{thm:6.5.2} by showing that $a^{(6,3^{n-2})}_{(1^n),(3)}=1$ for all $n\ge 3$ (note when $n\le 5$ this follows by direct computation).%\hfill$\lozenge$
\end{example}

A corollary of Theorem~\ref{thm:14.6i} is the stability of the following sequence of plethysm coefficients, whose monotonicity was predicted in \cite[Conjecture 1.2]{BBP}.

\begin{proposition}\label{prop:17}
	Let $\lambda$ and $\mu$ be partitions. For all $j\in\N_0$, define partitions $\lambda^j\coloneqq \lambda\sqcup(1^j)$ and $\mu^j\coloneqq (\mu+(j))\sqcup(1^j)$. Then the sequence $\big(a^{\mu^j}_{\lambda^j,(2)}\big)_{j\in\N_0}$ is eventually constant.
\end{proposition}

To prove Proposition~\ref{prop:17}, we first record a stability property of Littlewood--Richardson coefficients. For convenience we include a proof in our present notation.
\begin{lemma}\label{lem:lr-stab}
	Let $\alpha,\beta,\gamma$ and $\delta$ be partitions. Define $\gamma(j)\coloneqq \gamma+(j)$ and $\delta(j)\coloneqq \delta+(j)$ for all $j\in\N_0$. Then the sequence $\big( \langle\chi^{\gamma(j)/\alpha},\chi^{\delta(j)/\beta}\rangle \big)_{j\in\N_0}$ is non-decreasing and eventually constant.
\end{lemma}

\begin{proof}
	For a skew shape $\rho$, let $\cL(\rho)$ denote the set of all Littlewood--Richardson fillings of $\rho$. Let $\ff_j:\cL(\gamma(j)/\alpha)\to \cL(\gamma(j+1)/\alpha)$ be the map given by filling the additional box in $\gamma(j+1)/\alpha - \gamma(j)/\alpha$ with the number 1. Clearly $\ff_j$ is well-defined and injective for all $j$, and furthermore, bijective for all sufficiently large $j$ (for example, $j\ge \alpha_1+\gamma_2$ will suffice).
	Similarly define $\fg_j:\cL(\delta(j)/\beta)\to\cL(\delta(j+1)/\beta)$. For $\mathsf{s}\in\cL(\gamma(j)/\alpha)$ and $\mathsf{t}\in\cL(\delta(j)/\beta)$, note that $\mathsf{s}$ and $\mathsf{t}$ have the same type if and only if $\ff_j(\mathsf{s})$ and $\fg_j(\mathsf{t})$ have the same type. The assertion of the lemma follows since $\langle\chi^{\gamma(j)/\alpha},\chi^{\delta(j)/\beta}\rangle = |\{(\mathsf{s},\mathsf{t})\in\cL(\gamma(j)/\alpha)\times\cL(\delta(j)/\beta) \mid \mathsf{s}, \mathsf{t}\text{ have the same type} \}|$.
\end{proof}

\begin{proof}[Proof of Proposition~\ref{prop:17}]
	We may assume that $\lambda\vdash n$, $\mu\vdash 2n$ and $l(\mu)=n-k$ for some $n$ and $k\in\N_0$, else $a^{\mu^j}_{\lambda^j,(2)}=0$ for all $j$ by Lemma~\ref{lem:tall-pleth}. Since $\mu^j\vdash 2n+2j$ and $l(\mu^j)=(n+j)-k$, we have from Theorem~\ref{thm:14.6i} that
	\[ a^{\mu^j}_{\lambda^j,(2)} = \sum_{i=0}^k (-1)^{k+i}\cdot \sum_{\substack{\alpha\vdash k+i\\\beta\vdash i}} a^{\alpha/(k-i)}_{\beta',(2)}\cdot a^{\hat{\mu^j}/\alpha}_{{\lambda^j}'/\beta,(1)} \]
	for all $j$, where $\hat{\mu^j}\coloneqq \mu^j-(1^{n+j-k})$. The proof is concluded by observing that $a^{\hat{\mu^j}/\alpha}_{{\lambda^j}'/\beta,(1)}=\langle \chi^{{\lambda^j}'/\beta},\chi^{\hat{\mu^j}/\alpha}\rangle$ and using Lemma~\ref{lem:lr-stab} with $\gamma\coloneqq \hat{\mu}$ and $\delta\coloneqq \lambda'$.
\end{proof}

\begin{remark}
	A similar argument can be used to give a new proof that the sequence $\big(a^{\mu^j}_{\lambda^j,(2)}\big)_j$ also stabilises where $\lambda^j\coloneqq \lambda+(j)$ and $\mu^j=\mu\sqcup(2^j)$; this sequence is already known to be both non-decreasing and eventually constant by \cite[\textsection 2.6 Corollary 1]{Brion}.%\hfill$\lozenge$
\end{remark}



\subsection{Proof of Theorem~\ref{thm:14.6i}}\label{sec:16}
We first introduce some notation in preparation for the proofs to come.
\begin{notation}\label{not:rho}
	\begin{enumerate}[label=\textup{(\roman*)}]
		\item \label{561} Let $\lambda$ be a partition, $n\in\N_0$ and $\phi$ be a virtual character of $S_n$.
		We define
		\[ \phi/\chi^\lambda\coloneqq  \sum_{\mu\vdash n}\langle \phi,\chi^\mu\rangle\cdot\chi^{\mu/\lambda} \]
		where $\chi^{\mu/\lambda}=0$ if $\lambda\not\subseteq\mu$.
		
		\item \label{562} For $m,n\in\N$ and $\alpha/\beta$ a skew shape of size $n$, define \[
		\rho^{\alpha/\beta}_m \coloneqq  \cX(\triv_{S_m};\chi^{\alpha/\beta})\up_{S_m\wr S_n}^{S_{mn}}.\] 
		
		\item \label{563} For $\phi\in\Ch(S_{n_1})$ and $\theta\in\Ch(S_{n_2})$, define
		\[ \phi\boxtimes\theta\coloneqq  (\phi\times\theta)\up_{S_{n_1}\times S_{n_2}}^{S_{n_1+n_2}}. \]
		
		\item \label{564} Let $S_\lambda$ denote the Young subgroup $S_{\lambda_1}\times\cdots\times S_{\lambda_{l(\lambda)}}$ of $S_n$ and let
		\[ \zeta^\lambda\coloneqq \triv_{S_\lambda}\up^{S_n} = \chi^{(\lambda_1)}\boxtimes\cdots\boxtimes \chi^{(\lambda_{l(\lambda)})} \]
		denote the character of the permutation module $M^\lambda$. (Note $\zeta^\mu=\zeta^\lambda$ if $\mu$ is a composition of $n$ with the same parts as $\lambda$ but in a different order.) 
	\end{enumerate}
\end{notation}

Recall that the irreducible decomposition of such a permutation character $\zeta^\lambda$ is described by Young's Rule \cite[2.8.5]{JK}. Equivalently,
\begin{equation}\label{eqn:kostka}
	\zeta^\lambda = \sum_{\gamma\vdash n} K_{\gamma,\lambda}\cdot \chi^\gamma
\end{equation}
where the Kostka number $K_{\gamma,\lambda} = \langle \zeta^\lambda,\chi^\gamma\rangle=c^\gamma_{(\lambda_1),\dotsc,(\lambda_{l(\lambda)})}$ equals the number of semistandard Young tableaux of shape $\gamma$ and content $\lambda$.
In particular, we therefore have
\begin{equation}\label{eqn:15.2}
	\sum_{\gamma\vdash n} K_{\gamma,\lambda}\cdot \rho^\gamma_m = \cX(\triv_{S_m}; \zeta^\lambda)\up_{S_m\wr S_n}^{S_{mn}} = \rho^{(\lambda_1)}_m\boxtimes\cdots\boxtimes\rho^{(\lambda_{l(\lambda)})}_m.
\end{equation}

We now precisely identify the $i=k$ term of Theorem~\ref{thm:14.6i}, in Theorem~\ref{thm:14.6ii}. We then deduce Theorem~\ref{thm:14.6i} from Theorem~\ref{thm:14.6ii}. Following that, we prove Theorem~\ref{thm:14.6ii}, during the course of which we will also prove Theorem B, which has been numbered as Theorem~\ref{thm:16.7} in this section for ease of reference.

\begin{theorem}\label{thm:14.6ii}
	Fix $n\in\N$. Let $m\in\N$, $k\in\{0,1,\dotsc,n-1\}$ and $\lambda\vdash n$. Let $\nu\vdash mn+k$ with $l(\nu)=n$, and set $\hat{\nu}\coloneqq \nu-(1^n)\vdash (m-1)n+k$. Then
	\[ a^{\nu/(1^k)}_{\lambda',(m)} = \sum_{\substack{\alpha\vdash mk\\\beta\vdash k}} a^\alpha_{\beta',(m)}\cdot a^{\hat{\nu}/\alpha}_{\lambda/\beta,(m-1)}. \]
\end{theorem}

\begin{proof}[Proof of Theorem~\ref{thm:14.6i} from Theorem~\ref{thm:14.6ii}]
	We proceed by induction on $k$, observing that the $k=0$ case of Theorem~\ref{thm:14.6i} follows immediately from the $k=0$ case of Theorem~\ref{thm:14.6ii}. 
	Now assume $k>0$, and fix $\mu\vdash mn$ with $l(\mu)=n-k$. Let $\nu=\mu\sqcup(1^k)$, so $l(\nu)=n$ and $\hat{\nu}=\hat{\mu}$.
	
	We aim to evaluate $a^{\nu/(1^k)}_{\lambda',(m)}$, and compare it to Theorem~\ref{thm:14.6ii}. To do this, we will study the constituents in the skew character $\chi^{\nu/(1^k)}$. First, note that for any $\omega\vdash|\nu|-k$, by Theorem~\ref{thm:LR} we must have $c^\nu_{\omega,(1^k)}\in\{0,1\}$. Moreover, $c^\nu_{\omega,(1^k)}=1$ if and only if $\omega\subseteq\nu$ and all $k$ boxes of $[\nu/\omega]$ belong to different rows. We will denote by $\mathcal{A}$ the collection of $\omega$ with $c^\nu_{\omega,(1^k)}=1$. In particular,
	\[ \sum_{\omega\in\mathcal{A}}\chi^\omega = \chi^{{\nu}/(1^k)}. \]
	We partition $\mathcal{A}$ as a disjoint union, $\mathcal{A} = \coprod_{j=0}^{k}\mathcal{A}_j$, where $\mathcal{A}_j$ is the collection of $\omega\in\mathcal{A}$ for which $l(\omega)=n-k+j$. Notice that for each $j$, $\mathcal{A}_j$ bijects to $\mathcal{B}_j \coloneqq  \{\varpi\vdash|\hat{\mu}|-j \mid c^{\hat{\mu}}_{\varpi,(1^j)}\}$; the map is given by removal of the first column, $\omega\mapsto\hat{\omega}$, and this is seen to be a bijection by Theorem~\ref{thm:LR}. By a similar application of Theorem~\ref{thm:LR},
	
	\begin{equation}\label{eqn:B_j}
		\sum_{\varpi\in\mathcal{B}_j}\chi^\varpi = \chi^{\hat{\mu}/(1^j)}.
	\end{equation}
	Now observe that
	\[ a^{\nu/(1^k)}_{\lambda',(m)} = \sum_{\omega\in\mathcal{A}} a^{\omega}_{\lambda',(m)} = a^{\mu}_{\lambda',(m)} + \sum_{j=1}^k \sum_{\omega\in\mathcal{A}_j} a^{\omega}_{\lambda',(m)}. \]
	The idea here is that in our partition of $\mathcal{A}$, the $j=0$ term contributes precisely $a^{\mu}_{\lambda',(m)}$. Let us set $X \coloneqq  -\sum_{j=1}^k \sum_{\omega\in\mathcal{A}_j} a^{\omega}_{\lambda',(m)}$. Assuming Theorem~\ref{thm:14.6ii}, it suffices to show that $X$ equals the $\sum_{i=0}^{k-1}(\dotsc)$ part of the summation on the right hand side of \eqref{eqn:A}. For $0<j\leq k$, we have that
	\begin{align*}
		\sum_{\omega\in\mathcal{A}_j} a^{\omega}_{\lambda',(m)} &\overset{\text{\normalfont\tiny(inductive hypothesis)}}{=} \sum_{\omega\in\mathcal{A}_j} \sum_{i=0}^{k-j}(-1)^{k-j+i}\sum_{\substack{\gamma\vdash k-j + (m-1)i\\ \beta\vdash i}}a^{\gamma/(k-j-i)}_{\beta',(m)}\cdot a^{\hat{\omega}/\gamma}_{\lambda/\beta,(m-1)}\\
		&\overset{\text{\normalfont\tiny(bijection $\mathcal{A}_j\to\mathcal{B}_j$)}}{=}\sum_{\varpi\in\mathcal{B}_j} \sum_{i=0}^{k-j}(-1)^{k-j+i}\sum_{\substack{\gamma\vdash k-j + (m-1)i\\ \beta\vdash i}}a^{\gamma/(k-j-i)}_{\beta',(m)}\cdot a^{\varpi/\gamma}_{\lambda/\beta,(m-1)}
		\\
		&\overset{\normalfont\tiny \text{\eqref{eqn:a}}}{=}\sum_{\varpi\in\mathcal{B}_j} \sum_{i=0}^{k-j}(-1)^{k-j+i}\sum_{\substack{\gamma\vdash k-j + (m-1)i\\ \beta\vdash i}}a^{\gamma/(k-j-i)}_{\beta',(m)}\cdot \left\langle\chi^{\varpi},\rho^{\lambda/\beta}_{m-1}\boxtimes\chi^\gamma\right\rangle \\
		&\overset{\normalfont\tiny\text{\eqref{eqn:B_j}}}{=}\sum_{i=0}^{k-j}(-1)^{k-j+i}\sum_{\substack{\gamma\vdash k-j + (m-1)i\\ \beta\vdash i}}a^{\gamma/(k-j-i)}_{\beta',(m)}\cdot \left\langle\chi^{\hat{\mu}/(1^j)},\rho^{\lambda/\beta}_{m-1}\boxtimes\chi^\gamma\right\rangle
		\\
		&=\sum_{i=0}^{k-j}(-1)^{k-j+i}\sum_{\substack{\gamma\vdash k-j + (m-1)i\\ \beta\vdash i}} a^{\gamma/(k-j-i)}_{\beta',(m)}\cdot \left\langle\chi^{\hat{\mu}},\rho^{\lambda/\beta}_{m-1}\boxtimes\chi^\gamma\boxtimes\chi^{(1^j)}\right\rangle\\
		&=\sum_{i=0}^{k-j}(-1)^{k-j+i}\sum_{\substack{\gamma\vdash k-j + (m-1)i\\ \alpha\vdash k+(m-1)i\\ \beta\vdash i}} c^{\alpha}_{\gamma,(1^j)}\cdot a^{\gamma/(k-j-i)}_{\beta',(m)}\cdot \left\langle\chi^{\hat{\mu}},\rho^{\lambda/\beta}_{m-1}\boxtimes\chi^\alpha\right\rangle\\
		&\overset{\normalfont\tiny\text{\eqref{eqn:a}}}{=}\sum_{i=0}^{k-j}(-1)^{k-j+i}\sum_{\substack{\gamma\vdash k-j + (m-1)i\\ \alpha\vdash k+(m-1)i\\ \beta\vdash i}} c^{\alpha}_{\gamma,(1^j)}\cdot a^{\gamma/(k-j-i)}_{\beta',(m)}\cdot a^{\hat{\mu}/\alpha}_{\lambda/\beta,(m-1)}.
	\end{align*}
	It follows that
	\begin{align*}
		X &= -\sum_{j=1}^k \sum_{i=0}^{k-j}(-1)^{k-j+i}\sum_{\substack{\gamma\vdash k-j + (m-1)i\\ \alpha\vdash k+(m-1)i\\ \beta\vdash i}} c^{\alpha}_{\gamma,(1^j)}\cdot a^{\gamma/(k-j-i)}_{\beta',(m)}\cdot a^{\hat{\mu}/\alpha}_{\lambda/\beta,(m-1)}\\
		&= \sum_{i=0}^{k-1} \sum_{j=1}^{k-i} (-1)^{k+i} \cdot (-1)^{j+1} \sum_{\substack{\gamma\vdash k-j + (m-1)i\\ \alpha\vdash k+(m-1)i\\ \beta\vdash i}} c^{\alpha}_{\gamma,(1^j)}\cdot a^{\gamma/(k-j-i)}_{\beta',(m)}\cdot a^{\hat{\mu}/\alpha}_{\lambda/\beta,(m-1)}\\
		&= \sum_{i=0}^{k-1} (-1)^{k+i} \sum_{j=1}^{k-i} \sum_{\beta\vdash i} (-1)^{j+1}\cdot Y^\beta_{i,j},
	\end{align*}
	where we set
	\begin{align*}
		Y^\beta_{i,j}\coloneqq \sum_{\substack{\gamma\vdash k-j + (m-1)i\\ \alpha\vdash k+(m-1)i}} c^{\alpha}_{\gamma,(1^j)}\cdot a^{\gamma/(k-j-i)}_{\beta',(m)}\cdot a^{\hat{\mu}/\alpha}_{\lambda/\beta,(m-1)}.
	\end{align*}
	Next, we will simplify $Y^\beta_{i,j}$, for any fixed $\beta$, $i$ and $j$. To ease notation, in the rest of this proof we will abbreviate sums over all partitions of a given size. That is, we shorten $\sum_{\omega\vdash t}$ to $\sum_\omega$ (the size $t$ will always be clear from context). We use \eqref{eqn:skew-plethysm} to obtain
	\begin{align*}
		Y^\beta_{i,j}=\sum_{\gamma} \sum_\alpha c^\alpha_{\gamma,(1^j)}\cdot a^{\hat{\mu}/\alpha}_{\lambda/\beta,(m-1)} \cdot \sum_\varepsilon c^\gamma_{\varepsilon,(k-j-i)}\cdot a^\varepsilon_{\beta',(m)}.
	\end{align*}
	By \cite[(2.1)]{SLthesis}, we have that $\sum_\gamma c^\alpha_{\gamma,(1^j)}\cdot c^{\gamma}_{\varepsilon,(k-j-i)} = c^{\alpha}_{\varepsilon,(k-j-i),(1^j)} = \langle \chi^{\alpha/\varepsilon}, \chi^{k-j-i}\boxtimes \chi^{(1^j)}\rangle$. By Theorem~\ref{thm:LR}, $\chi^{k-j-i}\boxtimes \chi^{(1^j)} = \chi^{H(j)}+\chi^{H(j-1)}$ where $H(j)\coloneqq (k-i-j,1^j)$, except if $j=k-i$ then $\chi^\emptyset\boxtimes \chi^{(1^{k-i})}=\chi^{H(k-i-1)}$ only, i.e.~we treat the $\chi^{H(k-i)}$ term as the zero character. Hence
	\begin{align*}
		Y^\beta_{i,j} &= \sum_\alpha \sum_\varepsilon \langle \chi^{\alpha/\varepsilon}, \chi^{H(j)}+\chi^{H(j-1)}\rangle \cdot a^\varepsilon_{\beta',(m)}\cdot a^{\hat{\mu}/\alpha}_{\lambda/\beta,(m-1)}\\
		&= \sum_\alpha \sum_\varepsilon (c^\alpha_{\varepsilon,H(j)} + c^\alpha_{\varepsilon,H(j-1)}) \cdot a^\varepsilon_{\beta',(m)}\cdot a^{\hat{\mu}/\alpha}_{\lambda/\beta,(m-1)}
	\end{align*}
	(where we omit the $c^\alpha_{\varepsilon,H(k-i)}$ term). Since $H(0)=(k-i)$, we obtain
	\[ \sum_{j=1}^{k-i}(-1)^{j+1} \cdot Y^\beta_{i,j} = \sum_\alpha \sum_\varepsilon c^\alpha_{\varepsilon,(k-i)}\cdot a^\varepsilon_{\beta',(m)}\cdot a^{\hat{\mu}/\alpha}_{\lambda/\beta,(m-1)} = \sum_\alpha a^{\alpha/(k-i)}_{\beta',(m)}\cdot a^{\hat{\mu}/\alpha}_{\lambda/\beta,(m-1)}.\]
	Thus, we finally obtain
	\begin{align*} X &= \sum_{i=0}^{k-1}(-1)^{k+i} \cdot \sum_{\beta\vdash i} \sum_{j=1}^{k-i}(-1)^{j+1}\cdot Y^\beta_{i,j} \\&= \sum_{i=0}^{k-1}(-1)^{k+i}\cdot \sum_{\beta\vdash i}\ \sum_{\alpha\vdash k+(m-1)i} a^{\alpha/(k-i)}_{\beta',(m)}\cdot a^{\hat{\mu}/\alpha}_{\lambda/\beta,(m-1)}, \end{align*}
	as desired.
\end{proof}

To prove Theorem~\ref{thm:14.6ii}, we first deal with the case of $m=1$. In this case, $a^{\nu/(1^k)}_{\lambda',(1)} = c^{\nu}_{\lambda',(1^k)}$, which takes value 1 if $\nu-\lambda'$ is a sequence of 0s and 1s containing exactly $k$ many 1s, and takes value 0 otherwise. On the other hand, $\sum_{\alpha,\beta\vdash k} \big(a^\alpha_{\beta',(1)}\cdot a^{\hat{\nu}/\alpha}_{\lambda/\beta,\emptyset}\big)=c^{\lambda'}_{\hat{\nu},(1^{n-k})}$, which takes value 1 if $\lambda'-\hat{\nu}$ is a sequence of 0s and 1s containing exactly $n-k$ many 1s, and takes value 0 otherwise. We see that these two quantities are equal since $\nu=\hat{\nu}+(1^n)$, and hence Theorem~\ref{thm:14.6ii} holds when $m=1$. 

Next, we introduce some lemmas in preparation for proving Theorem~\ref{thm:14.6ii} when $m\ge 2$.

\begin{lemma}\label{lem:16.2}
	Let $r\in\N$ and $n_1,\dotsc,n_r\in\N_0$. Let $\gamma$ be a partition and for each $i\in\{1,2,\dotsc,r\}$, let $\psi_i$ be a virtual character of $S_{n_i}$. Then
	\[ (\psi_1\boxtimes\cdots\boxtimes\psi_r)/\chi^\gamma = \sum_{\gamma_1,\dotsc,\gamma_r} c^\gamma_{\gamma_1,\dotsc,\gamma_r}\cdot (\psi_1/\chi^{\gamma_1})\boxtimes\cdots\boxtimes(\psi_r/\chi^{\gamma_r}), \]
	summed over all sequences of partitions $\gamma_1,\dotsc,\gamma_r$ such that $|\gamma_1|+\cdots+|\gamma_r|=|\gamma|$.
\end{lemma}

\begin{proof}
	The case $r=1$ is trivial. For ease of notation we prove the statement for $r=2$; the case of general $r$ follows by an analogous argument. Define $n\coloneqq n_1+n_2$, $k\coloneqq |\gamma|$ and let $\delta\vdash n-k$. Then
	\begin{align*}
		\langle (\psi_1\boxtimes\psi_2)/\chi^\gamma,\chi^\delta\rangle &= \langle (\psi_1\times\psi_2)\up_{S_{n_1}\times S_{n_2}}^{S_n}, (\chi^\delta\times\chi^\gamma)\up_{S_{n-k}\times S_k}^{S_n}\rangle\\
		&= \langle (\psi_1\times\psi_2)\up_{S_{n_1}\times S_{n_2}}^{S_n}\down_{S_{n-k}\times S_k}, \chi^\delta\times\chi^\gamma\rangle,
	\end{align*}
	and applying Mackey's Theorem,
	\footnotesize
	\begin{align*}
		&= \sum_{\substack{0\le t_1,t_2\le k\\ t_1+t_2=k}} \langle (\psi_1\times\psi_2)\down^{S_{n_1}\times S_{n_2}}_{S_{n_1-t_1}\times S_{t_1}\times S_{n_2-t_2}\times S_{t_2}}\up^{S_{n-k}\times S_{k}}, \chi^\delta\times\chi^\gamma\rangle\\
		&= \sum_{t_1+t_2=k} \sum_{\substack{i=1,2\\\gamma_i\vdash t_i\\\delta_i\vdash n_i-t_i}} c^\gamma_{\gamma_1,\gamma_2}\cdot c^\delta_{\delta_1,\delta_2}\cdot \langle \psi_1\down_{S_{n_1-t_1}\times S_{t_1}} \times \psi_2\down_{S_{n_2-t_2}\times S_{t_2}}, \chi^{\delta_1}\times\chi^{\gamma_1}\times\chi^{\delta_2}\times\chi^{\gamma_2}\rangle\\
		&= \sum_{t_1+t_2=k} \sum_{\substack{i\\\gamma_i\vdash t_i\\\delta_i\vdash n_i-t_i}} c^\gamma_{\gamma_1,\gamma_2}\cdot c^\delta_{\delta_1,\delta_2}\cdot \langle \psi_1/\chi^{\gamma_1},\chi^{\delta_1}\rangle\cdot\langle \psi_2/\chi^{\gamma_2},\chi^{\delta_2}\rangle\\
		&= \sum_{t_1+t_2=k} \sum_{\substack{i\\\gamma_i\vdash t_i}} c^\gamma_{\gamma_1,\gamma_2}\cdot \langle (\psi_1/\chi^{\gamma_1})\boxtimes(\psi_2/\chi^{\gamma_2}), \chi^\delta\rangle,
	\end{align*}
	\normalsize
	as claimed.
\end{proof}

\begin{corollary}\label{cor:16.2}
	Let $m,n\in\N$ and $\lambda=(\lambda_1,\dotsc,\lambda_r)\vdash n$. Let $k\in\{0,1,\dotsc,n-1\}$ and $\beta\vdash k$. Then
	\begin{enumerate}[label=\textup{(\roman*)}]
		\item \label{591} $\big(\operatorname*{\boxtimes}_{i=1}^r \rho^{(\lambda_i)}_m \big)/\chi^{(1^{n-k})} = \sum_{\underline{t}}  \operatorname*{\boxtimes}_{i=1}^r \big( \rho^{(\lambda_i)}_m/\chi^{(1^{\lambda_i-t_i})} \big)$, and
		
		\item \label{592} $\big( \operatorname*{\boxtimes}_{i=1}^r \chi^{(1^{\lambda_i})} \big)/\chi^{\beta'} = \sum_{\underline{t}} c^\beta_{(t_1),\dotsc,(t_r)}\cdot \operatorname*{\boxtimes}_{i=1}^r \chi^{(1^{\lambda_i-t_i})}$,
	\end{enumerate}
	summed over all compositions $\underline{t}=(t_1,\dotsc,t_r)$ of $k$ into $r$ parts. That is, $t_i\in\N_0$ for all $i$ and $t_1+\cdots+t_r=k$ \textup{(}and we may further assume $t_i\le\lambda_i$ for all $i$\textup{)}.
\end{corollary}

\begin{proof}
	\ref{591} Applying Lemma~\ref{lem:16.2} with $\gamma=(1^{n-k})$, observe that $c^\gamma_{\gamma_1,\dotsc,\gamma_r}\in\{0,1\}$ and is non-zero only if each $\gamma_i=(1^{s_i})$ for some $s_i\in\N_0$. By Lemma~\ref{lem:tall-pleth}, $\rho^{(\lambda_i)}_m$ only has irreducible constituents $\chi^\mu$ where $l(\mu)\le \lambda_i$, so we may further assume that $s_i\le\lambda_i$. Writing $t_i=\lambda_i-s_i$ gives the result.
	
	\noindent \ref{592} Applying Lemma~\ref{lem:16.2} with $\psi_i=\chi^{(1^{\lambda_i})}$, observe that $\psi_i/\chi^{\gamma_i}\ne 0$ only if $\gamma_i=(1^{t_i})$ for some $0\le t_i\le \lambda_i$. Moreover, $c^{\beta'}_{(1^{t_1}),\dotsc,(1^{t_r})}=c^\beta_{(t_1),\dotsc,(t_r)}\in\{0,1\}$.
\end{proof}

\begin{lemma}\label{lem:16.3}
	Let $n\in\N$ and $k\in\{0,1,\dotsc,n\}$. Let $\nu$ be a partition with $l(\nu)=n$ and set $\hat{\nu}\coloneqq \nu-(1^n)$. Let $\delta\vdash|\nu|-k$ with $l(\delta)\le n$. Then
	\[ \langle\chi^\delta,\chi^{\nu/(1^k)} \rangle = \langle \chi^{\delta/(1^{n-k})},\chi^{\hat{\nu}}\rangle. \]
\end{lemma}

\begin{proof}
	The determinantal form of skew characters of symmetric groups (see e.g.~\cite[2.3.13]{JK}) gives $\chi^{\alpha/\beta} = \det\big(\chi^{(\alpha_i-i-\beta_j+j)}\big)$ whenever $\alpha$ and $\beta$ are partitions, where the multiplication of characters in expanding the determinant is given by the operation $\boxtimes$. Applying this to $\alpha=\nu'$ and $\beta=\emptyset$ and expanding the determinant with respect to the first row gives
	\[ \chi^{\nu'} = \sum_{j\ge 1} (-1)^{j-1}\cdot \chi^{(n+j-1)}\boxtimes \chi^{\hat{\nu}'/(1^{j-1})} = \sum_{j\ge 0} (-1)^j\cdot\chi^{(n+j)}\boxtimes \chi^{\hat{\nu}'/(1^j)}. \]
	Multiplying both sides by the sign representation then gives
	\[ \chi^\nu = \sum_{j\ge 0} (-1)^j\cdot \chi^{(1^{n+j})}\boxtimes \chi^{\hat{\nu}/(j)}. \]
	Then
	\begin{align*}
		\langle\chi^\delta,\chi^{\nu/(1^k)}\rangle &= \langle \chi^\delta\boxtimes\chi^{(1^k)},\chi^\nu\rangle \\
		&= \sum_{j\ge 0}(-1)^j\cdot\langle(\chi^\delta\boxtimes\chi^{(1^k)})/\chi^{(1^{n+j})}, \chi^{\hat{\nu}/(j)} \rangle\\
		&\overset{\normalfont\tiny \text{(Lemma~\ref{lem:16.2})}}{=} \sum_{j=0}^k (-1)^j\cdot \left\langle \sum_{s=0}^{k-j}(\chi^\delta/\chi^{(1^{n-s})})\boxtimes(\chi^{(1^k)}/\chi^{(1^{s+j})}), \chi^{\hat{\nu}/(j)}\right\rangle\\
		&=\sum_{j=0}^k \sum_{s=0}^{k-j}(-1)^j \langle (\chi^\delta/\chi^{(1^{n-s})})\boxtimes \chi^{(1^{k-j-s})}, \chi^{\hat{\nu}/(j)}\rangle\\
		&= \sum_{s=0}^k \left\langle (\chi^\delta/\chi^{(1^{n-s})})\boxtimes \left(\sum_{j=0}^{k-s}(-1)^j\cdot\chi^{(1^{k-s-j})}\boxtimes\chi^{(j)}\right), \chi^{\hat{\nu}}\right\rangle\\
		&=\langle \chi^\delta/\chi^{(1^{n-k})}, \chi^{\hat{\nu}}\rangle,
	\end{align*}
	where the final equality follows since $\chi^{(1^{k-s-j})}\boxtimes\chi^{(j)}=\chi^{(j+1,1^{k-s-j-1})}+\chi^{(j,1^{k-s-j})}$, and so $\sum_{j=0}^{k-s}(-1)^j\cdot\chi^{(1^{k-s-j})}\boxtimes\chi^{(j)}$ equals zero if $k\ne s$, and equals $\chi^\emptyset$ if $k=s$.
\end{proof}

\begin{lemma}\label{lem:16.6}
	Let $m$, $u$ and $t$ be integers with $m\ge 2$ and $u\ge t\ge 0$. Then
	\[ \rho^{(u)}_m/\chi^{(1^{u-t})} = \rho^{(t)}_m\boxtimes \rho^{(1^{u-t})}_{m-1}. \]
\end{lemma}

\begin{proof}
	Let $\delta\vdash mu-(u-t)$ be arbitrary. We show that $\langle \rho^{(u)}_m/\chi^{(1^{u-t})}, \chi^\delta\rangle = \langle\rho^{(t)}_m\boxtimes \rho^{(1^{u-t})}_{m-1}, \chi^\delta\rangle$. 
	Letting $H\coloneqq S_m\wr S_u$ and $K\coloneqq S_{|\delta|}\times S_{u-t}$, and substituting in the definition $\rho^{(u)}_m$ from Notation~\ref{not:rho}, we have by Mackey's theorem that
	\begin{align}\label{eqn:16.6}
		\langle \rho^{(u)}_m/\chi^{(1^{u-t})}, \chi^\delta\rangle &= \langle \rho^{(u)}_m, \chi^\delta\boxtimes\chi^{(1^{u-t})}\rangle = \langle \triv_H\up^{S_{mu}}, (\chi^\delta\times\chi^{(1^{u-t})})\up_K^{S_{mu}}\rangle \nonumber\\
		&= \sum_{\sigma\in K\setminus S_{mu}/H}\langle\triv_{H^\sigma}\down_{K\cap H^\sigma}\up^K,\chi^\delta\times\chi^{(1^{u-t})}\rangle\nonumber\\
		&=\sum_{\sigma\in K\setminus S_{mu}/H}\langle\triv_{K\cap H^\sigma},(\chi^\delta\times\chi^{(1^{u-t})})\down^K_{K\cap H^\sigma}\rangle,
	\end{align}
	where the final equality follows from Frobenius reciprocity.
	Here $\sigma$ runs over a set of representatives of double $(K,H)$-cosets in $S_{mu}$. Since $K\cap H^\sigma$ are the point stabilisers of the action of $K$ on the set of partitions of $\{1,2,\dotsc,mu\}$ into $u$ subsets of size $m$, the representatives $\sigma$ are parametrised by partitions of $u-t$ into exactly $u$ parts, including parts of size zero. 
	Fix one such partition $\sigma$ of $u-t$ and suppose that $\gamma_i$ is the number of parts of size $i$, for each $i\in\N_0$. Then $K\cap H^\sigma \cong \prod_{i\in\N_0} (S_{m-i}\times S_i)\wr S_{\gamma_i}$, and
	\begin{align*}
		\langle (\chi^\delta\times\chi^{(1^{u-t})})\down^K_{K\cap H^\sigma}, \triv\rangle &\le \langle (\chi^\delta\times\chi^{(1^{u-t})})\down_{\prod_i (S_{m-i}\times S_i)^{\times \gamma_i}}, \triv\rangle\\
		&= \langle \chi^\delta\down_{\prod_i S_{m-i}^{\times\gamma_i}},\triv\rangle\cdot \langle \chi^{(1^{u-t})}\down_{\prod_i S_i^{\times\gamma_i}},\triv\rangle.
	\end{align*}
	However, $\chi^{(1^{u-t})}$ is the sign representation, so $\langle \chi^{(1^{u-t})}\down_{\prod_i S_i^{\times\gamma_i}},\triv\rangle\ne 0$ if and only if $\gamma_i=0$ for all $i\ge 2$. Hence there is at most one $\sigma$ giving a non-zero contribution to the sum in \eqref{eqn:16.6}, namely $\sigma=(1^{u-t},0^t)$, and in this case $K\cap H^\sigma\cong (S_m\wr S_t)\times \big((S_{m-1}\times S_1)\wr S_{u-t}\big)$. Substituting into \eqref{eqn:16.6},
	\begin{align*}
		\langle \rho^{(u)}_m/\chi^{(1^{u-t})}, \chi^\delta\rangle &= \langle (\chi^\delta\times\chi^{(1^{u-t})})\down^K_{S_m\wr S_t\times (S_{m-1}\times S_1)\wr S_{u-t}},\ \triv\rangle,
		\end{align*}
which by Frobenius reciprocity equals
		\[\langle (\chi^\delta\times\chi^{(1^{u-t})})\down^{S_{|\delta|}\times S_{u-t}}_{S_m\wr S_t\times S_{m-1}\wr S_{u-t}\times S_1\wr S_{u-t}},\ \triv\up_{S_m\wr S_t\times (S_{m-1}\times S_1)\wr S_{u-t}}^{{S_m\wr S_t\times S_{m-1}\wr S_{u-t}\times S_1\wr S_{u-t}}}\rangle.
	\]
		 Noting that $|\delta|=mu-(u-t)=mt+(m-1)(u-t)$ and $\cX(\triv_{S_1};\chi^\omega)=\chi^\omega$, and using Lemma~\ref{lem:X-induced} in the second equality below, we have
	\begin{align*}
		\langle \rho^{(u)}_m/\chi^{(1^{u-t})}, \chi^\delta\rangle = \left\langle \chi^\delta\down^{S_{|\delta|}}_{S_m\wr S_t\times S_{m-1}\wr S_{u-t}} \times \chi^{(1^{u-t})},\ \triv_{S_m\wr S_t} \times \triv\up_{(S_{m-1}\times S_1)\wr S_{u-t}}^{S_{m-1}\wr S_{u-t}\times S_1\wr S_{u-t}}\right\rangle,
		\end{align*}
which in turn equals
\[
 \left\langle \chi^\delta\down^{S_{|\delta|}}_{S_m\wr S_t\times S_{m-1}\wr S_{u-t}} \times \chi^{(1^{u-t})},\ \triv_{S_m\wr S_t} \times \sum_{\omega\vdash u-t}\cX(\triv_{S_{m-1}}; \chi^\omega)\cdot\cX(\triv_{S_1};\chi^\omega) \right\rangle.
 \]
 This last expression simplifies to
\[
 \sum_{\omega\vdash u-t} \langle \chi^\delta\down_{S_m\wr S_t\times S_{m-1}\wr S_{u-t}},\ \triv_{S_m\wr S_t}\times \cX(\triv_{S_{m-1}};\chi^\omega)\rangle\cdot\langle \chi^{(1^{u-t})},\chi^\omega\rangle.
 \]
Now $\langle \chi^{(1^{u-t})},\chi^\omega\rangle=1$ precisely when $\omega=(1^{u-t})$ and is 0 otherwise, so
	\begin{align*} \langle \rho^{(u)}_m/\chi^{(1^{u-t})}, \chi^\delta\rangle &= \langle \chi^\delta\down_{S_m\wr S_t\times S_{m-1}\wr S_{u-t}},\ \triv_{S_m\wr S_t}\times \cX(\triv_{S_{m-1}};\chi^{(1^{u-t})})\rangle\\ &= \langle \chi^\delta,\ \rho^{(t)}_m\boxtimes \rho^{(1^{u-t})}_m\rangle \end{align*}
	by Frobenius reciprocity, recalling Notation~\ref{not:rho}\ref{562} and \ref{563}.
	Since $\delta$ was arbitrary, then $\rho^{(u)}_m/\chi^{(1^{u-t})} = \rho^{(t)}_m\boxtimes \rho^{(1^{u-t})}_{m-1}$ as desired.
\end{proof}

\begin{remark}
	When $m=2$, we can see from Proposition~\ref{prop:thrall} that $\rho^{(u)}_2/\chi^{(1^{u-t})} = \sum_\delta \chi^\delta = \rho^{(t)}_2\boxtimes\chi^{(1^{u-t})}$ where the sum is over all $\delta\vdash u+t$ with exactly $u-t$ many odd parts.%\hfill$\lozenge$
\end{remark}

Next, we generalise Lemma~\ref{lem:16.6} from the trivial partition $(u)$ to arbitrary partitions, giving Theorem B, after which it will be straightforward to deduce Theorem~\ref{thm:14.6ii}.

\begin{theorem}[Theorem B]\label{thm:16.7}
	Let $m,n\in\N$ with $m\ge 2$. Let $\lambda\vdash n$ and $k\in\{0,1,\dotsc,n-1\}$. Then
	\[ \rho^\lambda_m/\chi^{(1^{n-k})} = \sum_{\beta\vdash k} \rho^\beta_m\boxtimes \rho^{\lambda'/\beta'}_{m-1}. \]
\end{theorem}

\begin{proof}
	Following the notation for $\underline{t}$ in Corollary~\ref{cor:16.2} and letting $r=l(\lambda)$, observe that
	\begin{align*}
		\sum_{\gamma\vdash n} K_{\gamma,\lambda} &\cdot \big( \rho^\gamma_m/\chi^{(1^{n-k})} \big) = \Big( \sum_{\gamma\vdash n} K_{\gamma,\lambda}\cdot \rho^\gamma_m \Big)/\chi^{(1^{n-k})} \overset{\normalfont\tiny\text{\eqref{eqn:15.2}}}{=} \Big( \operatorname*{\boxtimes}_{i=1}^r \rho^{(\lambda_i)}_m\Big)/\chi^{(1^{n-k})}\\
		&\overset{\normalfont\tiny\text{(Corollary~\ref{cor:16.2}\ref{591})}}{=} \sum_{\underline{t}} \operatorname*{\boxtimes}_{i=1}^r \big( \rho^{(\lambda_i)}_m/\chi^{(1^{\lambda_i-t_i})} \big)\\
		&\overset{\normalfont\tiny\text{Lemma~\ref{lem:16.6}}}{=} \sum_{\underline{t}} \operatorname*{\boxtimes}_{i=1}^r \big( \rho^{(t_i)}_m \boxtimes \rho^{(1^{\lambda_i-t_i})}_{m-1} \big)\\
		&\overset{\normalfont\tiny\text{\eqref{eqn:15.2}}}{=} \sum_{\underline{t}} \sum_{\beta\vdash k} c^\beta_{(t_1),\dotsc,(t_r)}\cdot\rho^\beta_m \boxtimes \big(\operatorname*{\boxtimes}_{i=1}^r \rho^{(1^{\lambda_i-t_i})}_{m-1}\big)\\
		&\overset{\normalfont\tiny \text{Lemma~\ref{lem:cX}}}{=} \sum_{\beta\vdash k} \rho^\beta_m \boxtimes \cX\Big(\triv_{S_{m-1}}; \sum_{\underline{t}} c^\beta_{(t_1),\dotsc,(t_r)}\cdot \operatorname*{\boxtimes}_{i=1}^r \chi^{(1^{\lambda_i-t_i})} \Big)\up_{S_{m-1}\wr S_{n-k}}^{S_{(m-1)(n-k)}}\\
		&\overset{\normalfont\tiny\text{Corollary~\ref{cor:16.2}\ref{592}}}{=} \sum_{\beta\vdash k} \rho^\beta_m \boxtimes \cX \Big( \triv_{S_{m-1}}; \big( \operatorname*{\boxtimes}_{i=1}^r \chi^{(1^{\lambda_i})}\big) /\chi^{\beta'} \Big) \up_{S_{m-1}\wr S_{n-k}}^{S_{(m-1)(n-k)}}\\
		&\overset{\normalfont\tiny\text{\eqref{eqn:sign}}}{=} \sum_{\beta\vdash k} \rho^\beta_m \boxtimes \cX \Big( \triv_{S_{m-1}}; \big( \zeta^\lambda\cdot\sgn_{S_n} \big) /\chi^{\beta'} \Big) \up_{S_{m-1}\wr S_{n-k}}^{S_{(m-1)(n-k)}}\\
		&\overset{\normalfont\tiny\text{\eqref{eqn:sign} and \eqref{eqn:kostka}}}{=} \sum_{\beta\vdash k} \rho^\beta_m \boxtimes \cX \Big( \triv_{S_{m-1}}; \big( \sum_{\gamma\vdash n} K_{\gamma,\lambda}\cdot \chi^{\gamma'} \big) /\chi^{\beta'} \Big) \up_{S_{m-1}\wr S_{n-k}}^{S_{(m-1)(n-k)}} \\
		&= \sum_{\gamma\vdash n} K_{\gamma,\lambda} \cdot \Big( \sum_{\beta\vdash k} \rho^\beta_m \boxtimes \rho^{\gamma'/\beta'}_{m-1} \Big)\quad\text{since $\chi^{\gamma'}/\chi^{\beta'}=\chi^{\gamma'/\beta'}$ by Notation~\ref{not:rho}\ref{561}}.
	\end{align*}
	Since the matrix $(K_{\gamma,\lambda})_{\gamma,\lambda\vdash n}$ is invertible (in fact unitriangular if the partitions are ordered lexicographically, see e.g.~\cite[Chapter 2]{JK}), we deduce that $\rho^\gamma_m/\chi^{(1^{n-k})} = \sum_{\beta\vdash k} \rho^\beta_m \boxtimes \rho^{\gamma'/\beta'}_{m-1}$ for each $\gamma\vdash n$.
\end{proof}

\begin{lemma}\label{lem:cX}
	Let $m,r,a_1,\dotsc,a_r\in\N$, and let $n=\sum_{i=1}^r a_i$. For each $i\in\{1,\dotsc,r\}$, let $\nu_i\vdash a_i$. Then $\boxtimes_{i=1}^r \rho_m^{\nu_i} = \cX(\triv_{S_m};\boxtimes_{i=1}^r \chi^{\nu_i} )\up_{S_m\wr S_n}^{S_{mn}}$.
\end{lemma}

\begin{proof}
	The case $r=1$ follows from Notation~\ref{not:rho}\ref{562}. For each of notation we prove the statement for $r=2$; the case of general $r$ follows by an analogous argument. In fact, we can prove more generally that if $a,b\in\N$ and $\phi_1\in\Ch(S_a), \phi_2\in\Ch(S_b)$, then $\cX_{left}=\cX_{right}$ where
	\[ \cX_{left}\coloneqq  \left[\cX(\triv_{S_m}; \phi_1)\up_{S_m\wr S_a}^{S_{ma}} \times \cX(\triv_{S_m}; \phi_2)\up_{S_m\wr S_b}^{S_{mb}}\right]\up_{S_{ma}\times S_{mb}}^{S_{m(a+b)}} \]
	and
	\[ \cX_{right}\coloneqq  \cX\left(\triv_{S_m}; (\phi_1\times \phi_2)\up_{S_a\times S_b}^{S_{a+b}}\right)\up_{S_m\wr S_{a+b}}^{S_{m(a+b)}}, \]
	from which we recover the case of $r=2$ by setting $\phi_i=\chi^{\nu_i}$. 
	
	To prove that $\cX_{left}=\cX_{right}$, we first observe that $S_m\wr(S_a\times S_b) = S_m\wr S_a \times S_m\wr S_b$ (viewing $S_a\times S_b$ as a subgroup of $S_{a+b}$). Calling this group $U$, it is a subgroup of both $T_d\coloneqq S_{ma}\times S_{mb}$ and $T_w\coloneqq S_m\wr S_{a+b}$, and both $T_d$ and $T_w$ are subgroups of $S\coloneqq S_{m(a+b)}$. Now, by Lemma~\ref{lem:infl-ind} and the definition of $\cX(-;-)$,
	\begin{multline*}
		\cX(\triv_{S_m}; (\phi_1 \times \phi_2)\up_{S_a\times S_b}^{S_{a+b}}) = \cX(\triv_{S_m}; \phi_1\times\phi_2)\up_U^{T_w} = \left(\Infl_{S_a\times S_b}^{T_w}(\phi_1\times\phi_2)\right)\up_U^{T_w}\\
		= \left( \Infl_{S_a}^{S_m\wr S_a}(\phi_1)\times \Infl_{S_b}^{S_m\wr S_b}(\phi_2) \right)\up_U^{T_w} = \Big( \cX(\triv_{S_m};\phi_1)\times \cX(\triv_{S_m};\phi_2) \Big)\up_U^{T_w}.
	\end{multline*}
	Therefore
	\begin{multline*}
	 	\cX_{right} = \Big( \cX(\triv_{S_m};\phi_1)\times \cX(\triv_{S_m};\phi_2) \Big)\up_U^{T_w} \up_{T_w}^S \\
	 	= \Big( \cX(\triv_{S_m};\phi_1)\times \cX(\triv_{S_m};\phi_2) \Big)\up_U^{T_d} \up_{T_d}^S = \cX_{left},
	 \end{multline*}
	where the second equality follows from the transitivity of induction.
\end{proof}

\begin{proof}[Proof of Theorem~\ref{thm:14.6ii} when $m\ge 2$]
	Take $\langle -, \chi^{\hat{\nu}}\rangle$ in Theorem~\ref{thm:16.7} to obtain
	\[ \langle \rho^\lambda_m/\chi^{(1^{n-k})}, \chi^{\hat{\nu}} \rangle = \left\langle \sum_{\beta\vdash k} \rho^\beta_m\boxtimes \rho^{\lambda'/\beta'}_{m-1}, \chi^{\hat{\nu}} \right\rangle. \]
	By Lemma~\ref{lem:16.3}, $\langle \rho^\lambda_m/\chi^{(1^{n-k})}, \chi^{\hat{\nu}} \rangle = \langle \rho^\lambda_m, \chi^{\nu/(1^k)}\rangle = a^{\nu/(1^k)}_{\lambda,(m)}$. On the other hand,
	\[ \left\langle \sum_{\beta\vdash k} \rho^\beta_m\boxtimes \rho^{\lambda'/\beta'}_{m-1}, \chi^{\hat{\nu}} \right\rangle = \sum_{\beta\vdash k} \sum_{\alpha\vdash mk} \langle \rho^\beta_m,\chi^\alpha\rangle \cdot \langle \rho^{\lambda'/\beta'}_{m-1}, \chi^{\hat{\nu}/\alpha}\rangle = \sum_{\substack{\alpha\vdash mk\\\beta\vdash k}} a^\alpha_{\beta,(m)}\cdot a^{\hat{\nu}/\alpha}_{\lambda'/\beta',(m-1)}, \]
	which concludes the proof.
\end{proof}
